%%%%%%%%%%%%%%%%%%%%%%%%%%%%%%%%%%%%%%%%%%%%%%%%%%%%%%%%%%%%%%%%%%%%%%%%%%%%%%%%
%%
\cleardoubleoddpage%  Make sure to start each chapter on a new odd page
\chapter{Dataset Details}
\label{app:dataset}

% TODO: full demographic tables; visit-interval histograms; per-channel
% statistics.


\cleardoubleoddpage
\chapter{Extended Derivations}
\label{app:derivations}

% TODO: Monge--Amp\`ere equation derivation; monitor function for the
% multi-channel case; pushforward stability argument.


\cleardoubleoddpage
\chapter{Full Hyperparameters}
\label{app:hyperparameters}

% TODO: everything not in the main text - all architecture widths, all
% training schedules, random seeds.


\cleardoubleoddpage
\chapter{Additional Rollout Figures}
\label{app:rollouts}

% TODO: extra qualitative examples, including failure cases.


\cleardoubleoddpage
\chapter{Ablation Details}
\label{app:ablations}

% TODO: full tables for every ablation sweep, not just the headline numbers.


\cleardoubleoddpage
\chapter{Code Structure}
\label{app:code}

% TODO: brief map of the repository, pointer to GitHub, reproducibility
% instructions.


\cleardoubleoddpage
\chapter{The Moving-Mesh Extension (MM-PDE)}
\label{app:mm-pde}

% Moved here from Chapter 2 (former \S2.5) on 2026-09-26, following the
% author's decision of 2026-09-25 to report everything MM-PDE as an
% additional experiment (THESIS_STRUCTURE.md, Appendix G).

\section{Background: MM-PDE}
\label{app:mm-pde:background}

\citet{Hu2024} extend message-passing neural solvers by integrating an
adaptive moving mesh. A uniform mesh is fundamentally inefficient for
problems that exhibit localised dynamics. It wastes computational
resolution in smooth regions while simultaneously under-resolving sharp
gradients, such as wave fronts.

Mesh adaptation strategies separate into two families. The
$h$-adaptation approach adds and removes vertices to dynamically refine
or coarsen the mesh. This changes both the topology and the total
degrees of freedom between steps, a property that is hostile to
standard graph-network training. Conversely, $r$-adaptation, or
moving-mesh adaptation, moves a fixed set of nodes while maintaining
both the node count and the original graph topology. Prior neural
$r$-adaptation methods were entirely supervised and required
ground-truth optimal meshes generated by an expensive classical solver.

The moving-mesh framework is governed by a monitor function. The
monitor is defined as $M = (1 + \lVert \nabla u \rVert_2 / \alpha) I$,
a matrix-valued function that takes large values precisely where the
state solution varies sharply. The identity factor $I$ ensures that the
refinement is isotropic. The addition of a constant $+1$ keeps $M$
bounded away from zero, which prevents nodes from collapsing entirely
in smooth regions. The global scaling constant $\alpha$ dictates the
trade-off between mesh adaptivity and numerical stability. The
corresponding mesh density is given by $\rho(x) = \sqrt{\det M(x)}$.
This specific functional form is derived from an established upper
bound on finite-element interpolation error.

The equidistribution principle formalises the objective of the mesh
adaptation. A coordinate map $f$ produces an $M$-uniform mesh if the
integral of the density $\rho$ over the image of every single cell is
equal. Specifically, this integral must equal $\sigma / N$, where
$\sigma$ is the total integral of $\rho$ over the domain and $N$ is the
total number of cells. Intuitively, every cell must carry the identical
amount of monitor-weighted information, which forces cells to shrink in
regions where $M$ is large.

Because infinitely many maps satisfy the equidistribution principle,
choosing the optimal map is framed as an optimal transport problem. The
objective is to select the mapping closest to the identity. By
Brenier's theorem, the unique solution to this problem is the gradient
of a convex scalar potential, such that $f = \nabla \phi$. Substituting
this gradient into the equidistribution constraint yields the
Monge--Amp\`ere equation,
\[
  \det H(\phi) \;=\; \frac{\sigma}{|\Omega| \, \rho(\nabla \phi)},
\]
which is a fully non-linear second-order PDE for the scalar potential
$\phi$. The problem is constrained by the boundary condition that the
domain boundary maps exactly to itself. A scalar potential is learned
rather than the full vector map because the convexity of $\phi$ is
mathematically equivalent to a non-singular Jacobian of $f$. A
non-singular Jacobian guarantees that the resulting mesh does not
tangle. Crucially, enforcing convexity on a scalar potential is far
easier to regularise than imposing a global non-tangling constraint on
a vector field.

The network architecture reformulates the scalar potential residually
as $\phi(x) = 0.5 \|x\|^2 + \psi(x)$. The gradient therefore becomes
$\nabla \phi = x + \nabla \psi$. The network is tasked with predicting
only the displacement field $\psi$. This formulation removes the
trivial identity component from the network's objective and ensures
that the coordinate map is approximately the identity mapping at
initialisation.

The Data-free Mesh Mover (DMM) follows a DeepONet-style architecture. A
branch network encodes the full spatial state $u$, and a trunk network
processes a continuous query coordinate. Together, they output the
scalar potential. The DMM is trained entirely without mesh supervision
via a dedicated physics loss
$L = L_{\text{eq}} + \beta L_{\text{bound}} + \gamma L_{\text{convex}}$.
In this formulation, $L_{\text{eq}}$ is the squared residual of the
Monge--Amp\`ere equation evaluated at sampled interior points,
$L_{\text{bound}}$ enforces the boundary condition, and
$L_{\text{convex}}$ penalises violations of convexity. The convexity
term is critical; a non-convex potential inevitably yields a tangled
and useless mesh, regardless of how small the equation residual
becomes. Training uses the Adam optimiser followed by L-BFGS
fine-tuning applied exclusively to the last layer. The DMM is
pretrained and subsequently frozen during solver training. This
decouples a well-posed geometric problem from the main solver objective
and allows the identical mesh mover to be reused across different
experiments.

The full moving-mesh update operator is constructed as a dual-branch
composition, defined as $M(u_t) = G_1(u_t) + I(G_2, \tilde{f}, u_t)$.
Here, $G_1$ represents a graph neural network operating on the uniform
mesh, $G_2$ is a graph neural network operating on the moved mesh,
$\tilde{f}$ is the coordinate map predicted by the DMM, and $I$ is a
learned interpolation operator mapping features back to uniform
coordinates. The uniform branch supplies a stable global view of the
domain, while the moved branch supplies a locally refined correction
where gradients are steep.

The interpolation network, ItpNet, learns adaptive interpolation
weights directly from the moved-mesh geometry rather than relying on
fixed finite-element basis functions. It is pretrained with its own
optimiser before joint dual-branch training begins, because an
untrained interpolation would feed unstructured noise into the uniform
branch, destabilising the system.

The classical background for adaptive moving meshes on which this
framework builds is covered extensively by \citet{HuangRussell2011}.
\citet{Hu2024} reported several architectural ablations that confirm
the dual-branch necessity. Removing the moved branch degrades
performance entirely to the uniform-mesh baseline. Removing the uniform
branch degrades performance even further, because the moved mesh
inherently under-covers smooth regions. Replacing the DMM-generated
mesh with a random mesh isolates and degrades the performance of the
moved branch specifically. Finally, removing the interpolation
pretraining severely harms overall training stability, even if it does
not permanently cap the final accuracy.

Whether mesh adaptation provides a measurable benefit for slowly
evolving clinical data is the question this appendix tests.

% TODO: remaining Appendix G content (THESIS_STRUCTURE.md): the DMM for GA
% (native-grid scalar monitor, physics loss; former \S4.5), the dual-branch
% composition with the alpha gate, the result (dual - bypass =
% -0.0001 +- 0.0060, scoped as the correction branch's failure to
% optimise), the mesh-quality study (former \S5.5) and the ~10x cost.


%%
%%%%%%%%%%%%%%%%%%%%%%%%%%%%%%%%%%%%%%%%%%%%%%%%%%%%%%%%%%%%%%%%%%%%%%%%%%%%%%%%
